\documentclass[11pt]{article}
\usepackage[margin=1in]{geometry}
\usepackage{amsmath, amssymb}
\usepackage{enumitem}
\usepackage{titlesec}
\usepackage{hyperref}
\usepackage{booktabs}
\usepackage{xcolor}
\usepackage{parskip}

\hypersetup{colorlinks=true, linkcolor=blue, urlcolor=blue}

\title{\textbf{EE627 --- Class Summary}\\[0.4em]
       \large January 23, 2026}
\author{Prof.\ Rensheng Wang}
\date{}

\begin{document}
\maketitle
\tableofcontents
\newpage

%------------------------------------------------------------
\section{Course Philosophy: What This Class Is Actually For}
%------------------------------------------------------------

\subsection{The Core Goal}
This is not a class about memorizing commands. The objective is to build the kind of analytical judgment that separates junior-level workers from senior-level ones --- and that separates candidates who get screened out from candidates who get hired.

The professor's framing: \textit{anything covered here, you can find on YouTube. The value is in how it is explained --- plain English, sharp points, with direct links to what industry actually tests.}

\subsection{The AI Shift Changes What Is Valuable}
\begin{itemize}
    \item ChatGPT, Gemini, and Copilot can generate code more precisely than most humans.
    \item Therefore, \textbf{writing code is no longer your competitive edge.}
    \item What AI \emph{cannot} provide is:
    \begin{itemize}
        \item Domain knowledge and context
        \item Judgment about whether a result is good or bad
        \item The decision to improve rather than accept
        \item Strategic direction: what features to engineer, which model to pick, how to iterate
    \end{itemize}
    \item The new interview standard: give candidates data and tools, then ask \textit{``Can you improve it?''} --- not \textit{``Can you write the syntax?''}
\end{itemize}

\textit{Analogy from lecture:} In the near future, every real engineer manages a team of AI agents. Each agent handles implementation. \textbf{You supply the judgment that tells the agents what to do.}

\subsection{Feature Engineering as the Differentiator}
The single most important skill emphasized throughout the semester is \textbf{feature engineering} --- deriving new, informative columns from raw data.

\begin{itemize}
    \item Raw data fed directly into a model is like a chef throwing an unpeeled potato into a wok. The preparation \emph{is} the craft.
    \item Feature engineering is open-ended: there is no one right answer, which is exactly why it reveals the depth of a candidate's thinking.
    \item Pre-screening tests at companies like Two Sigma and Goldman Sachs are frequently feature engineering problems for this reason.
\end{itemize}

\textit{``Squeeze the orange''} --- the job is to extract as much useful signal (``juice'') from the raw data as possible. How you do that is your insight.

\newpage
%------------------------------------------------------------
\section{Course Structure and Logistics}
%------------------------------------------------------------

\subsection{Assignments}
\begin{itemize}
    \item \textbf{Weekly homework} --- builds grades incrementally; each assignment gives multiple days to complete.
    \item \textbf{Midterm} --- approximately the first week of April (week of March 30 -- April 6).
    \item \textbf{Final project} --- end of semester.
    \item No single high-stakes test; weekly work compounds into the final grade.
\end{itemize}

\subsection{Submission Format}
\begin{itemize}
    \item Accepted: \textbf{PDF}, images (JPEG/PNG screenshots), or \texttt{.py} Python files.
    \item \textbf{Not accepted:} Jupyter notebook \texttt{.ipynb} files directly --- export to PDF first, or take a screenshot.
    \item \textbf{Not accepted:} Zip archives. Upload individual files only.
\end{itemize}

\textit{Why:} Canvas cannot render \texttt{.ipynb} files online. PDF and images allow the professor to annotate directly with brush/pencil comments, which is how grading feedback is delivered.

\subsection{Kaggle Platform}
Homework and projects are submitted through \textbf{Kaggle}:
\begin{itemize}
    \item Kaggle provides \textbf{instant scoring feedback} --- no waiting for the professor to grade to know if you are on the right track.
    \item All activity is logged: number of submissions, approaches tried, data used. There is no way to claim effort that is not recorded.
    \item Kaggle's open discussion forums let you see how peers (mostly beginners, like you) approach the same problems.
\end{itemize}

\newpage
%------------------------------------------------------------
\section{Setup Checklist for Week 2}
%------------------------------------------------------------

Before the next session, complete all of the following:

\begin{enumerate}
    \item \textbf{ChatGPT or Gemini account} --- free tier is sufficient. Use it to generate code from descriptions, then evaluate whether the output makes sense.
    \item \textbf{Python environment} with \texttt{xgboost} installed. Verify you can import the library.
    \item \textbf{Kaggle account} --- required to submit homework and receive scores.
    \item \textbf{GitHub familiarity (optional but recommended)} --- learn to browse open-source repos to avoid reinventing solutions to low-level problems.
\end{enumerate}

\textit{The professor's message:} Do not spend a weekend figuring out how to open a file in Python. That is a one-line ChatGPT question. Spend your learning budget on judgment, not syntax.

\newpage
%------------------------------------------------------------
\section{Course Topics Overview (Semester Syllabus)}
%------------------------------------------------------------

\subsection{Feature Engineering + XGBoost (Weeks 1--3)}
\begin{itemize}
    \item \textbf{XGBoost is the default first model for structured data.} If you have tabular data and no strong prior, start here.
    \item Structured data = anything with named columns (CSV, TSV, PSV). Every column has a defined meaning.
    \item After running XGBoost: if performance is mediocre, the next question is always \textit{``How do I improve it?''} --- that question leads directly to feature engineering.
\end{itemize}

\subsection{Time Series Analysis}
\begin{itemize}
    \item Time series = data indexed in a fixed sequential order that cannot be shuffled.
    \item Core concept: \textbf{autocorrelation} --- exploiting the dependency between adjacent time steps to forecast future values.
    \item Broader relevance: the same index-dependency logic appears in:
    \begin{itemize}
        \item Natural language (word order in a sentence)
        \item Geospatial data (neighboring locations)
        \item Generative AI diffusion models
        \item Transformers
    \end{itemize}
    \item Financial markets are one of the richest sources of time series data, and a strong reason to develop this skill set.
\end{itemize}

\subsection{Linear and Logistic Regression (Conceptual Bridge)}
\begin{itemize}
    \item The goal is not to learn \emph{how} to run these models (you already know that).
    \item The goal is to understand \emph{why they fail} in certain settings --- which directly motivates the jump to neural networks.
    \item Understanding limitations gives you language to explain poor model performance and propose the right fix.
\end{itemize}

\newpage
\subsection{Recommendation Systems (Main Focus)}
The semester's central application domain. Reasons:
\begin{enumerate}
    \item Recommendation is a superset of feature engineering, big data, and model improvement --- all major course themes appear here.
    \item It is one of the highest cash-flow areas in tech: a 0.001\% improvement to a system used by billions of users translates to massive profit. This explains why top engineers in recommendation/search earn dramatically more than industry average.
    \item Nobody can escape it: Google, Amazon, Netflix, Instagram, TikTok --- every platform's output is filtered through recommendation.
\end{enumerate}

\textit{Recommendation and search are two sides of the same coin.} Flip a recommendation engine around and you get a search engine. The underlying methods are nearly identical.

\subsection{Ensemble Methods}
\begin{itemize}
    \item Instead of relying on one model, combine multiple models' outputs --- different judgments fused into one stronger decision.
    \item Covered toward the end of the semester as a capstone technique.
\end{itemize}

\newpage
%------------------------------------------------------------
\section{Structured vs.\ Unstructured Data}
%------------------------------------------------------------

\begin{center}
\begin{tabular}{@{}lll@{}}
\toprule
\textbf{Type}        & \textbf{Examples}                        & \textbf{Default Model}     \\
\midrule
Structured           & CSV, TSV, SQL tables                     & XGBoost                    \\
Unstructured (image) & Pixel arrays, photos                     & Deep learning / CNN        \\
Unstructured (text)  & Articles, sentences, chat logs           & Deep learning / Transformer \\
\bottomrule
\end{tabular}
\end{center}

\textbf{Why XGBoost does not work on images or text:}
A pixel at position (row 1, col 1) has no fixed semantic meaning --- it could represent an eye, a background, noise, anything. Structured data's power comes from every column having a stable, interpretable meaning. Without that, tree-based methods have nothing meaningful to split on.

Unstructured data requires neural networks (covered in EE628) to learn representations automatically.


\newpage
%------------------------------------------------------------
\section{Industry Context: How AI Changes the Job Market}
%------------------------------------------------------------

\begin{itemize}
    \item Pure front-end and back-end software roles are increasingly automatable.
    \item ML/AI engineering is the growth area --- and it requires humans for the judgment layer.
    \item The new interview format: open data + open tools. The question is not \textit{``Write a function''} but \textit{``Here is a dataset. Improve the model.''} Show multiple approaches. That range of thinking is what impresses interviewers.
    \item A candidate who says \textit{``I ran random forest''} and stops has shown a ceiling. A candidate who proposes four different feature engineering directions demonstrates potential.
\end{itemize}

\textbf{The multiplier effect:} A strong engineer who uses AI as a force multiplier produces output equivalent to 10--100 weaker engineers. The key word is \emph{multiplier} --- it amplifies what you bring in, not a substitute for it.

\end{document}
